% Preamble: \usepackage{tikz}  \usetikzlibrary{arrows.meta}
\begin{figure}[t]
\centering
\resizebox{\columnwidth}{!}{%
\begin{tikzpicture}[
    font=\footnotesize,
    >={Latex[length=1.6mm]},
    card/.style={draw=violet!70!black, fill=violet!10, rounded corners=2pt,
                 align=center, inner sep=1.5pt, minimum width=24mm,
                 minimum height=4.2mm, font=\scriptsize},
    out/.style={draw=#1!70!black, fill=#1!10, rounded corners=2pt,
                align=center, inner sep=2pt, minimum width=46mm,
                minimum height=7mm, font=\scriptsize}
]

% ---- Top: established k=1 ----
\node[font=\scriptsize\bfseries] at (0, 3.05) {Established ($k{=}1$): insight once};
\node[card, draw=black!45, fill=black!6] (q1) at (-1.6, 2.45) {Problem $q_i$};
\node[card] (i1) at (1.6, 2.45) {Insight $e_i$};
\node[out=teal, minimum width=52mm] (base) at (0, 1.65)
    {accuracy preserved, fewer generated tokens};

% ---- Middle: the intervention ----
\node[font=\scriptsize\bfseries] at (0, 0.95) {This paper ($k>1$): repeat the \emph{same} insight};
\node[card, draw=black!45, fill=black!6] (q2) at (-1.6, 0.35) {Problem $q_i$};
\node[card] (ia) at (1.6, 0.50) {Insight $e_i$};
\node[card] (ib) at (1.6, 0.05) {Insight $e_i \;(k\times)$};

% ---- Bottom: three outcomes ----
\node[font=\scriptsize\bfseries] at (0, -0.55) {Which outcome?};
\node[out=teal,    minimum width=52mm] (o1) at (0, -1.15)
    {\textbf{Help:} reinforces the hint};
\node[out=black!30, fill=black!6, minimum width=52mm] (o2) at (0, -1.85)
    {\textbf{Saturate:} extra copies add nothing};
\node[out=orange,  minimum width=52mm] (o3) at (0, -2.55)
    {\textbf{Hurt:} over-constrains, inflates prompt};

% ---- Arrows ----
\draw[->] (0, 1.35) -- (0, 0.80);
\draw[->] (0, -0.10) -- (0, -0.40);
\node[font=\Large\bfseries, text=black!70] at (2.55, -0.25) {?};

% ---- Footnote ----
\node[font=\scriptsize, text=black!55] at (0, -3.15)
    {Only $k$ changes; insight, question, model, decoding fixed.};
\end{tikzpicture}}
\caption{The question this paper isolates. An insight inserted \emph{once}
($k{=}1$) already preserves accuracy while cutting tokens. We repeat the
\emph{same} insight $k$ times, changing nothing else, and ask which of three
outcomes follows. \textbf{Takeaway:} the literature predicts all three, so the
sign is genuinely open---which is why it warrants a controlled test.}
\label{fig:teaser}
\end{figure}